Uma mesma faixa pode ser descrita pelo som, pelo contexto de lançamento e pelo conteúdo da letra. Essas fontes de informação podem ser parcialmente redundantes. Gênero e atributos acústicos, por exemplo, já descrevem aspectos do estilo musical que também aparecem no texto. A questão deste trabalho é se conhecer a letra melhora a classificação do nível de popularidade observado no Spotify quando áudio e metadados já estão disponíveis.

\citeonline{dhanaraj2005} investigaram a combinação de texto e áudio na classificação de sucesso musical. \citeonline{singhi2014} estudaram características de letras para distinguir músicas populares. Já \citeonline{parada2024lyrics} analisaram a contribuição textual em dados de escuta da Last.fm. A comparação direta dos desempenhos exige atenção às respostas de cada estudo. Neste trabalho, a resposta é o escore fornecido pelo Spotify no arquivo disponível. Para a análise, ele é dividido em três níveis ordenados.

O objetivo é medir a contribuição incremental de dois tipos de representação das letras, descritores lexicais e estruturais e embeddings multilíngues. A comparação usa quatro configurações ajustadas nas mesmas faixas e divisões, com o mesmo modelo estatístico. Examina-se tanto o desempenho médio entre classes quanto o tipo de erro que passa a ocorrer com cada representação. O desempenho refere-se à classificação do escore observado, sem um horizonte futuro de previsão.

A integração segue quatro passos. Primeiro, descrevem-se o catálogo acessível do Spotify e a estrutura da base LRCLIB. Depois, dimensiona-se e sorteia-se uma amostra probabilística do catálogo para orientar a recuperação e aferir sua cobertura. Em seguida, os vínculos aceitos dessa busca e de buscas anteriores são reunidos em um quadro com letras. Dele saem as 9.526 faixas usadas pelos quatro classificadores. A distinção entre o sorteio inicial e a seleção no quadro vinculado determina quais resultados podem ser atribuídos ao catálogo acessível e quais descrevem a coorte de modelagem.

